\section{Methodology}\label{sec:methodology}

\begin{figure}[!t]
\centering
\includegraphics[width=\linewidth]{analysis/figures/h_methodology.pdf}
\caption{\textbf{Methodology overview.} \textbf{(a)}~\emph{Harness refinement within one episode}: the Agent reads $(s_t,\mathcal{H},\tau)$ and emits $a_t$; every $F$ steps the Refiner reads $\tau_{t-F:t}$, emits per-component edits $\Delta=(\Delta p,\Delta\mathcal{G},\Delta\mathcal{K},\Delta\mathcal{M})$ via the meta-tool API, and $\mathcal{H}\!\gets\!\mathcal{H}\oplus\Delta$. \textbf{(b)}~\emph{Co-learning across DAgger+PRM iterations}: each iteration runs $\pi_{\theta_k}$ inside a live-refining $\mathcal{H}_t$ for $K{=}256$ steps. The trajectory is scored by a pairwise PRM, low-$R$ windows are relabeled by Gemini-3.1-pro, and a soft SFT update produces $\theta_{k+1}$. The loop is reset-free: a persistent state at the end of iter~$k$ is loaded as the start of iter~$k{+}1$.}
\label{fig:methodology}
\end{figure}

\subsection{Overview and Two-Loop Architecture}
\texttt{Continual Harness} performs online in-context learning over the harness state $\mathcal{H}$ from \cref{sec:preliminaries}. An LLM Refiner edits $\mathcal{H}$ from the most recent trajectory window during a single continuous episode, generalizing prompt-optimization methods that rewrite only $p$ from complete-episode resets~\citep{agrawal2025gepa, opsahl2024optimizing} to a method that rewrites the full state from the trajectory so far.

Write $s_t = (o_t, m_t)$ for the agent's observation at step $t$. The \emph{inner loop} is the standard agent step: the model $M$ wrapped by the current harness $\mathcal{H}_t$ produces an action $a_t$ from $s_t$ and the trajectory so far. The \emph{outer loop} is harness refinement: every $F$ steps after a warm-up of $W$ steps, a Refiner reads the recent trajectory window for failure signatures and emits per-component edits $\Delta = (\Delta p, \Delta\mathcal{G}, \Delta\mathcal{K}, \Delta\mathcal{M})$. The agent does not reset; the updated harness $\mathcal{H}_{t+1} = \mathcal{H}_t \oplus \Delta$ enters the agent's context on the next step (\cref{fig:methodology}a), with $p$ replaced by $\Delta p$ and $\mathcal{G}, \mathcal{K}, \mathcal{M}$ receiving CRUD-style operations (create, read, update, delete). The Agent and Refiner roles share the same model $M$, ablated across Gemini 3.1 Pro, Flash, and Flash-Lite (\cref{sec:experiments}). In our GPP runs, the Refiner role for the system prompt and pre-built primitives was performed manually by humans observing the livestream; \texttt{Continual Harness} automates it. Both the agent and the Refiner issue edits through the same meta-tool API (\cref{sec:preliminaries}); they differ only in when each is invoked and on what trajectory context.

\subsection{Refinement Loop}\label{sec:refinement}

The Refiner reads $\tau_{t-F:t}$ and identifies failure signatures over the window: navigation loops, tool-call failures, stalled objectives, and missed exploration opportunities. It then runs four passes, one per component: (i) it rewrites the prompt $p$ conditioned on the identified failures and the trajectory window; (ii) it creates sub-agent entries for repeated multi-step patterns, edits existing entries to address detected failures, and deletes entries that have not been invoked productively; (iii) it codifies skills from successful sequences and repairs executable code that raised exceptions; (iv) it adds memory entries to fill gaps, updates stale entries, and demotes importance for areas the agent has moved past.

Refinement information accumulates monotonically over the episode: failure signatures observed earlier in the trajectory remain available to all subsequent refinement passes, so refinement quality compounds with episode length, while reset-based methods restart this accumulation after each update. \texttt{Continual Harness} can also target failure modes that only appear deep in an episode (late-game battles, multi-step puzzles, dialogue chains), which reset-based approaches cannot reach by construction since each iteration resets to the initial state. Beyond these technical advantages, reset-free is also the practically dominant regime for long-running coding agents, embodied agents, and ops tasks where free environment resets are costly or unavailable.

\subsection{Continual Model-Harness Co-Learning Loop}\label{sec:training}

\cref{fig:methodology}b instantiates Continual Harness as a training loop for an open-source model. After warm-up stages (Appendix~\ref{app:training}), each online iteration runs $\pi_{\theta_k}$ inside a live-refining harness $\mathcal{H}_t$ for $K{=}256$ steps. A pairwise process reward model (PRM) $R(s_t, a_t, \tau) \in [0,1]$ scores each transition over a sliding window of recent transitions (component weights in Appendix~\ref{app:training}); low-reward windows are relabeled by a frontier teacher, and a soft SFT update on the relabeled shard produces $\theta_{k+1}$. The loop is reset-free since the saved emulator state at the end of iteration $k$ is loaded as the start of iteration $k{+}1$, so the model's in-game position accumulates across training rather than restarting.

The trajectory distribution $\mathcal{D}_\theta$ depends on $\theta$ through the harness. The model's actions induce $\tau$, the Refiner reads $\tau$ to update $\mathcal{H}_t$, and $\mathcal{H}_t$ in turn shapes the next observation distribution. Both the model weights $\theta$ and the harness state $\mathcal{H}_t$ are updated by this loop, where $\theta$ is updated across iterations (via SFT on relabeled trajectories) and $\mathcal{H}_t$ within each iteration (via the Refiner).
